\documentclass[11pt]{article}
\usepackage{ap-lit}
\title{Unit 1: Prose Fiction: Question Set}
\author{Paul Wang}
\date{November 2026}
\docpart{Unit 1 \textperiodcentered\ Question Set}

\begin{document}
\maketitle

\begin{abstract}
Two exam-style multiple-choice sets on public-domain fiction, followed by a full Prose Fiction Analysis essay. Passage 1 is the opening of James Joyce's ``Araby'' (1914), a retrospective first-person narration; Passage 2 is a short chapter from Kate Chopin's \emph{The Awakening} (1899), with an intrusive third-person narrator. Distractors follow the exam's patterns: half-right pairs, tones that are too extreme, and choices that are true of the passage but do not answer the question. Allow about 25 minutes for the multiple choice (about 65 seconds per question, including reading) and 40 minutes for the essay. Paragraph numbers are printed in the margin.
\end{abstract}

\scoringline{MCQ: \blank\ / \qtotal \qquad Essay: \blank\ / 6 \qquad Target: 17+ MCQ, 4+ essay}

\section{Passage 1: James Joyce, ``Araby''}

\begin{passage}{Questions 1--13.\quad Read the passage carefully before you choose your answers.}{The following passage is the opening of ``Araby,'' a short story by James Joyce published in 1914 in the collection \emph{Dubliners}.}
\pp North Richmond Street, being blind, was a quiet street except at the hour when the Christian Brothers' School set the boys free. An uninhabited house of two storeys stood at the blind end, detached from its neighbours in a square ground. The other houses of the street, conscious of decent lives within them, gazed at one another with brown imperturbable faces.

\pp The former tenant of our house, a priest, had died in the back drawing-room. Air, musty from having been long enclosed, hung in all the rooms, and the waste room behind the kitchen was littered with old useless papers. Among these I found a few paper-covered books, the pages of which were curled and damp: \emph{The Abbot}, by Walter Scott, \emph{The Devout Communicant} and \emph{The Memoirs of Vidocq}. I liked the last best because its leaves were yellow. The wild garden behind the house contained a central apple-tree and a few straggling bushes under one of which I found the late tenant's rusty bicycle-pump. He had been a very charitable priest; in his will he had left all his money to institutions and the furniture of his house to his sister.

\pp When the short days of winter came dusk fell before we had well eaten our dinners. When we met in the street the houses had grown sombre. The space of sky above us was the colour of ever-changing violet and towards it the lamps of the street lifted their feeble lanterns. The cold air stung us and we played till our bodies glowed. Our shouts echoed in the silent street. The career of our play brought us through the dark muddy lanes behind the houses where we ran the gauntlet of the rough tribes from the cottages, to the back doors of the dark dripping gardens where odours arose from the ashpits, to the dark odorous stables where a coachman smoothed and combed the horse or shook music from the buckled harness. When we returned to the street light from the kitchen windows had filled the areas. If my uncle was seen turning the corner we hid in the shadow until we had seen him safely housed. Or if Mangan's sister came out on the doorstep to call her brother in to his tea we watched her from our shadow peer up and down the street. We waited to see whether she would remain or go in and, if she remained, we left our shadow and walked up to Mangan's steps resignedly. She was waiting for us, her figure defined by the light from the half-opened door. Her brother always teased her before he obeyed and I stood by the railings looking at her. Her dress swung as she moved her body and the soft rope of her hair tossed from side to side.

\pp Every morning I lay on the floor in the front parlour watching her door. The blind was pulled down to within an inch of the sash so that I could not be seen. When she came out on the doorstep my heart leaped. I ran to the hall, seized my books and followed her. I kept her brown figure always in my eye and, when we came near the point at which our ways diverged, I quickened my pace and passed her. This happened morning after morning. I had never spoken to her, except for a few casual words, and yet her name was like a summons to all my foolish blood.

\pp Her image accompanied me even in places the most hostile to romance. On Saturday evenings when my aunt went marketing I had to go to carry some of the parcels. We walked through the flaring streets, jostled by drunken men and bargaining women, amid the curses of labourers, the shrill litanies of shop-boys who stood on guard by the barrels of pigs' cheeks, the nasal chanting of street-singers, who sang a come-all-you about O'Donovan Rossa, or a ballad about the troubles in our native land. These noises converged in a single sensation of life for me: I imagined that I bore my chalice safely through a throng of foes. Her name sprang to my lips at moments in strange prayers and praises which I myself did not understand. My eyes were often full of tears (I could not tell why) and at times a flood from my heart seemed to pour itself out into my bosom. I thought little of the future. I did not know whether I would ever speak to her or not or, if I spoke to her, how I could tell her of my confused adoration. But my body was like a harp and her words and gestures were like fingers running upon the wires.
\end{passage}

\begin{q}{2}{1.4 Retrospective narrator}{Past tense, ``foolish blood,'' ``confused adoration'': an adult recalling, and judging, his boyhood.}
The narrator of the passage is best described as
\ch{a boy reporting events as they happen to him}{an adult recalling his boyhood infatuation and at moments judging it}{an omniscient narrator who moves among the children's thoughts}{a neighbor observing the boy from across the street}
\end{q}

\begin{q}{4}{1.5 Connotation}{A ``blind'' street is a dead end; the word also suggests a life with no outlet or vision.}
In the first sentence of paragraph 1, the word ``blind'' most nearly suggests a street that
\ch{is too dark to see at night}{is unknown to visitors}{is dangerous to the children}{has no way through, hinting at lives that lead nowhere}
\end{q}

\begin{q}{1}{1.2 Atmosphere}{Musty air, useless papers, damp books, a rusty pump: decay left behind by the dead priest.}
The details in paragraph 2 (the musty air, the ``old useless papers,'' the damp books, the rusty bicycle-pump) chiefly create an atmosphere of
\ch{decay and neglect left behind by the past}{scholarly richness that inspires the boy}{comfortable domestic order}{suspense about a crime}
\end{q}

\begin{q}{3}{1.1 Indirect characterization}{He prefers a book for the color of its pages: a romantic, surface-drawn imagination.}
The narrator's comment that he liked \emph{The Memoirs of Vidocq} best ``because its leaves were yellow'' (paragraph 2) suggests that as a boy he
\ch{was a careful student of history}{disliked the priest who owned the books}{was drawn to things by their aura rather than their substance}{knew the books were valuable}
\end{q}

\begin{q}{3}{1.6 Irony}{``Charitable'' sits against money to institutions and only furniture to his sister.}
The last sentence of paragraph 2 (``He had been a very charitable priest \ldots'') is best read as
\ch{sincere praise that the narrator expects readers to share}{a foreshadowing of the boy's own career in the church}{quietly ironic, since the priest's ``charity'' went to institutions while his sister received only furniture}{a direct quotation of the priest's will}
\end{q}

\begin{q}{1}{1.5 Syntax}{One long sentence of ``through \ldots\ to \ldots\ to'' carries the reader along the boys' route.}
The long sentence beginning ``The career of our play'' (paragraph 3) is constructed chiefly to
\ch{carry the reader through a sequence of places, imitating the boys' run through the lanes}{slow the narrative to a single motionless image}{introduce the rough tribes as main characters}{contrast the boys' play with their schoolwork}
\end{q}

\begin{q}{4}{1.5 Imagery}{From ``our shadow'' the boys see her ``defined by the light'': an illuminated, idealized figure.}
In paragraph 3, the boys watch Mangan's sister ``from our shadow'' while her figure is ``defined by the light from the half-opened door.'' The contrast chiefly serves to
\ch{show that the boys are afraid of the dark}{reveal that Mangan's sister is unaware of the boys}{establish that the story takes place in summer}{present her as an illuminated, almost idealized figure seen from the dark}
\end{q}

\begin{q}{2}{1.1 Character}{Hiding behind the blind, following, passing her: devoted, secretive, inexperienced.}
In paragraph 4, the boy's actions (lying on the floor to watch her door with the blind ``pulled down to within an inch,'' following her, and passing her) characterize him as
\ch{confident and openly courting her}{devoted but secretive and unable to approach her}{indifferent to her except as a neighbor}{resentful of her attention to others}
\end{q}

\begin{q}{1}{1.5 Word choice}{``Never spoken \ldots\ except for a few casual words'' set against ``summons'' measures the gap.}
The sentence ``I had never spoken to her, except for a few casual words, and yet her name was like a summons to all my foolish blood'' (paragraph 4) chiefly emphasizes
\ch{the gap between the slightness of the relationship and the force of his feeling}{that the girl has rejected him}{that the narrator regrets speaking to her}{the narrator's skill as a speaker}
\end{q}

\begin{q}{3}{1.6 Allusion}{Chalice, litanies, prayers and praises: his love imagined as a sacred quest.}
In paragraph 5, the statement ``I imagined that I bore my chalice safely through a throng of foes,'' together with words such as ``litanies'' and ``prayers and praises,'' presents the boy's love as
\ch{a duty imposed on him by his aunt}{a joke he shares with the shop-boys}{a sacred quest, with the girl as an object of worship}{a political cause like the ballads he hears}
\end{q}

\begin{q}{2}{1.2 Setting}{The coarse, noisy market is ``hostile to romance''; his devotion moves through it untouched.}
The description of the Saturday markets in paragraph 5 chiefly serves to
\ch{show that the boy enjoys shopping with his aunt}{set the coarse, noisy city against the boy's private, idealized devotion}{introduce characters who will matter later in the story}{explain the economic causes of the troubles in Ireland}
\end{q}

\begin{q}{4}{1.6 Simile}{He is the harp and she the fingers: he is played upon, passive.}
The final sentence of the passage (``my body was like a harp \ldots'') suggests that the boy experiences his love as something
\ch{he has carefully planned}{that bores him}{shared equally by the girl}{that plays upon him from outside his control}
\end{q}

\begin{q}{1}{1.7 Literary argument}{Lyrical detail plus ``foolish,'' ``confused,'' ``could not tell why'': tender and gently ironic.}
Which claim about the passage as a whole is best supported?
\ch{The adult narrator regards his boyhood devotion with both tenderness and gentle irony.}{The narrator regrets that he never married Mangan's sister.}{The boy's love is openly returned by the girl.}{The neighborhood is presented as a place that encourages romance.}
\end{q}

\section{Passage 2: Kate Chopin, \emph{The Awakening}}

\begin{passage}{Questions 14--24.\quad Read the passage carefully before you choose your answers.}{The following passage is a complete chapter (Chapter VI) of \emph{The Awakening}, a novel by Kate Chopin published in 1899. Edna Pontellier, twenty-eight and married, is spending the summer at a resort on the Louisiana coast, where a young man, Robert, has asked her to go to the beach with him.}
\pp Edna Pontellier could not have told why, wishing to go to the beach with Robert, she should in the first place have declined, and in the second place have followed in obedience to one of the two contradictory impulses which impelled her.

\pp A certain light was beginning to dawn dimly within her,---the light which, showing the way, forbids it.

\pp At that early period it served but to bewilder her. It moved her to dreams, to thoughtfulness, to the shadowy anguish which had overcome her the midnight when she had abandoned herself to tears.

\pp In short, Mrs.\ Pontellier was beginning to realize her position in the universe as a human being, and to recognize her relations as an individual to the world within and about her. This may seem like a ponderous weight of wisdom to descend upon the soul of a young woman of twenty-eight---perhaps more wisdom than the Holy Ghost is usually pleased to vouchsafe to any woman.

\pp But the beginning of things, of a world especially, is necessarily vague, tangled, chaotic, and exceedingly disturbing. How few of us ever emerge from such beginning! How many souls perish in its tumult!

\pp The voice of the sea is seductive; never ceasing, whispering, clamoring, murmuring, inviting the soul to wander for a spell in abysses of solitude; to lose itself in mazes of inward contemplation.

\pp The voice of the sea speaks to the soul. The touch of the sea is sensuous, enfolding the body in its soft, close embrace.
\end{passage}

\begin{q}{3}{1.4 Omniscient narrator}{The narrator names the two impulses Edna ``could not have told'': it knows more than she can say.}
In paragraph 1, the narrator's statement that Edna ``could not have told why'' she acted as she did indicates that the narrator
\ch{is Edna herself, speaking years later}{knows nothing of Edna's thoughts}{understands Edna's conflict more clearly than Edna can}{doubts that Edna went to the beach at all}
\end{q}

\begin{q}{2}{1.6 Paradox}{The light reveals a path and, in revealing it, shows that it is forbidden.}
The description of ``the light which, showing the way, forbids it'' (paragraph 2) suggests that Edna's new awareness
\ch{is a religious conversion that ends her doubts}{reveals a possible path while also making clear that the path is prohibited}{is the sunrise she sees from the beach}{is a lesson Robert teaches her}
\end{q}

\begin{q}{4}{1.3 Structure}{Short, separated paragraphs isolate each stage of the awakening and give it weight.}
The chapter's division into very short paragraphs, several of a single sentence, chiefly serves to
\ch{speed up a sequence of physical actions}{reproduce a conversation between Edna and Robert}{list the events of a single day}{isolate each stage of an emerging awareness and give it weight}
\end{q}

\begin{q}{1}{1.6 Verbal irony}{``Usually pleased to vouchsafe to any woman'' mocks the assumption that women are not given such wisdom.}
The remark that this is ``perhaps more wisdom than the Holy Ghost is usually pleased to vouchsafe to any woman'' (paragraph 4) is best described as
\ch{wry irony aimed at conventional assumptions about women}{sincere religious reverence}{the narrator's contempt for Edna's foolishness}{a neutral report of Edna's education}
\end{q}

\begin{q}{2}{1.3 Shift}{``In short'' turns from figurative description to the narrator's direct summary.}
The phrase ``In short'' at the start of paragraph 4 marks a shift from
\ch{dialogue to description}{figurative description of Edna's state to the narrator's plain summary of it}{the present to a memory of Edna's childhood}{Edna's perspective to Robert's}
\end{q}

\begin{q}{3}{1.4 Narrator's generalization}{``How few of us'' widens Edna's experience to everyone's and raises the stakes.}
In paragraph 5, the exclamations ``How few of us ever emerge from such beginning! How many souls perish in its tumult!'' chiefly serve to
\ch{reveal that Edna will drown later that day}{quote Edna's words to Robert}{generalize from Edna's experience to human experience, raising its stakes}{mock readers who have never had such an experience}
\end{q}

\begin{q}{4}{1.5 Diction}{The piling adjectives mimic the disorder they name.}
The series ``vague, tangled, chaotic, and exceedingly disturbing'' (paragraph 5) is effective chiefly because it
\ch{softens the description into comfort}{describes the weather at the resort}{contradicts the narrator's earlier claims}{accumulates difficulty, so the sentence enacts the disorder it describes}
\end{q}

\begin{q}{2}{1.6 Personification}{A seductive voice ``inviting the soul'': the sea as tempter toward solitude.}
In paragraph 6, the personification of the sea's ``voice'' as ``seductive \ldots\ inviting the soul to wander'' presents the sea as
\ch{a threat of storms that Edna must avoid}{a tempter calling Edna toward solitude and inward life}{a symbol of her husband's authority}{a source of comic relief}
\end{q}

\begin{q}{1}{1.3 Shift}{Paragraph 6 turns from abstract reflection to sensuous description of the sea.}
The movement from paragraph 5 to paragraph 6 is best described as a shift from
\ch{abstract reflection on beginnings to sensuous description of the sea}{a description of Robert to a description of Edna}{dialogue to stage directions}{the present tense to a flashback}
\end{q}

\begin{q}{3}{1.5 Imagery}{The sea speaks to the soul, then touches the body: an awakening both spiritual and physical.}
The final paragraph, in which the sea first ``speaks to the soul'' and then enfolds ``the body in its soft, close embrace,'' suggests that Edna's awakening is
\ch{purely intellectual}{a memory she wishes to forget}{both spiritual and physical}{entirely imagined by the narrator}
\end{q}

\begin{q}{4}{1.7 Literary argument}{``Perish,'' ``abysses,'' ``mazes'' beside ``realize her position'': momentous and dangerous.}
Which claim about the chapter is best supported by the text?
\ch{The narrator disapproves of Edna's growing thoughtfulness.}{Edna fully understands her position by the end of the chapter.}{The sea represents the safety of domestic life.}{The narrator presents Edna's awakening as both momentous and dangerous.}
\end{q}

\section{Free Response: Prose Fiction Analysis}

\begin{frq}{Question 2 \textperiodcentered\ Prose Fiction Analysis \textperiodcentered\ 6 points \textperiodcentered\ 40 minutes}
The following excerpt is from ``Miss Brill,'' a short story by Katherine Mansfield first published in 1920. Miss Brill, an Englishwoman who earns her living in France as a teacher and reader, is spending her customary Sunday afternoon in a public garden, where a band plays. Read the passage carefully. Then, in a well-written essay, analyze how Mansfield uses literary elements and techniques to convey the complexity of Miss Brill's sense of belonging.

In your response you should do the following:
\begin{itemize}[leftmargin=1.4em]
\item Respond to the prompt with a thesis that presents a defensible interpretation.
\item Select and use evidence to support your line of reasoning.
\item Explain how the evidence supports your line of reasoning.
\item Use appropriate grammar and punctuation in communicating your argument.
\end{itemize}
\end{frq}

\begin{passage}{Katherine Mansfield, ``Miss Brill'' (1920), excerpt}{}
\pp Oh, how fascinating it was! How she enjoyed it! How she loved sitting here, watching it all! It was like a play. It was exactly like a play. Who could believe the sky at the back wasn't painted? But it wasn't till a little brown dog trotted on solemn and then slowly trotted off, like a little ``theatre'' dog, a little dog that had been drugged, that Miss Brill discovered what it was that made it so exciting. They were all on the stage. They weren't only the audience, not only looking on; they were acting. Even she had a part and came every Sunday. No doubt somebody would have noticed if she hadn't been there; she was part of the performance after all. How strange she'd never thought of it like that before! And yet it explained why she made such a point of starting from home at just the same time each week---so as not to be late for the performance---and it also explained why she had quite a queer, shy feeling at telling her English pupils how she spent her Sunday afternoons. No wonder! Miss Brill nearly laughed out loud. She was on the stage. She thought of the old invalid gentleman to whom she read the newspaper four afternoons a week while he slept in the garden. She had got quite used to the frail head on the cotton pillow, the hollowed eyes, the open mouth and the high pinched nose. If he'd been dead she mightn't have noticed for weeks; she wouldn't have minded. But suddenly he knew he was having the paper read to him by an actress! ``An actress!'' The old head lifted; two points of light quivered in the old eyes. ``An actress---are ye?'' And Miss Brill smoothed the newspaper as though it were the manuscript of her part and said gently; ``Yes, I have been an actress for a long time.''

\pp The band had been having a rest. Now they started again. And what they played was warm, sunny, yet there was just a faint chill---a something, what was it?---not sadness---no, not sadness---a something that made you want to sing. The tune lifted, lifted, the light shone; and it seemed to Miss Brill that in another moment all of them, all the whole company, would begin singing. The young ones, the laughing ones who were moving together, they would begin, and the men's voices, very resolute and brave, would join them. And then she too, she too, and the others on the benches---they would come in with a kind of accompaniment---something low, that scarcely rose or fell, something so beautiful---moving\ldots\ And Miss Brill's eyes filled with tears and she looked smiling at all the other members of the company. Yes, we understand, we understand, she thought---though what they understood she didn't know.

\pp Just at that moment a boy and girl came and sat down where the old couple had been. They were beautifully dressed; they were in love. The hero and heroine, of course, just arrived from his father's yacht. And still soundlessly singing, still with that trembling smile, Miss Brill prepared to listen.

\pp ``No, not now,'' said the girl. ``Not here, I can't.''

\pp ``But why? Because of that stupid old thing at the end there?'' asked the boy. ``Why does she come here at all---who wants her? Why doesn't she keep her silly old mug at home?''

\pp ``It's her fu-ur which is so funny,'' giggled the girl. ``It's exactly like a fried whiting.''

\pp ``Ah, be off with you!'' said the boy in an angry whisper. Then: ``Tell me, ma petite ch\`ere---''

\pp ``No, not here,'' said the girl. ``Not yet.''

\pp On her way home she usually bought a slice of honey-cake at the baker's. It was her Sunday treat. Sometimes there was an almond in her slice, sometimes not. It made a great difference. If there was an almond it was like carrying home a tiny present---a surprise---something that might very well not have been there. She hurried on the almond Sundays and struck the match for the kettle in quite a dashing way.

\pp But to-day she passed the baker's by, climbed the stairs, went into the little dark room---her room like a cupboard---and sat down on the red eiderdown. She sat there for a long time. The box that the fur came out of was on the bed. She unclasped the necklet quickly; quickly, without looking, laid it inside. But when she put the lid on she thought she heard something crying.
\end{passage}

\makekey

\subsection*{Essay scoring guide (6 points)}
\small
\begin{compare}{Row}{Earns the point(s)}{Does not earn the point(s)}
A. Thesis (1) & A defensible interpretation of the \emph{complexity} of her belonging, e.g., that the sense of belonging Miss Brill builds from watching others is real to her but rests on an illusion that one overheard remark is enough to destroy; or that she belongs to the scene only as a spectator who has cast herself as a performer. & ``Mansfield uses imagery, dialogue, and setting to show Miss Brill's sense of belonging''; ``Miss Brill is lonely''; a summary of the ending. \\
B. Evidence and commentary (4) & \textbf{4}: specific references from across the passage (the play metaphor and ``Even she had a part''; the invalid gentleman, whom she would not miss if he died; the music's ``faint chill''; ``we understand \ldots\ though what they understood she didn't know''; the couple as ``hero and heroine''; ``fried whiting''; the honey-cake almond; the ``room like a cupboard''; the fur ``crying'') with commentary that consistently explains how narration, imagery, structure, and dialogue build the interpretation. \textbf{3}: the same evidence, but only some of it is explained. \textbf{2}: some specific evidence with uneven, general, or partly mistaken commentary. \textbf{1}: summary of the story with little explanation. & Quoting the couple's insults and saying ``this is mean''; retelling the plot paragraph by paragraph; naming ``imagery'' and ``metaphor'' without explaining what they do. \\
C. Sophistication (1) & Develops a tension through the whole essay, e.g., how free indirect narration lets readers share Miss Brill's joy while seeing its self-deception; that the ``performance'' she imagines is true in a way she does not intend (she \emph{is} acting, for herself); or how the displaced ``crying'' shows a grief she still cannot claim as her own. & A final sentence about how ``everyone wants to belong''; a single flourish without development. \\
\end{compare}

\begin{modelessay}[A 6-point response, in outline]
\textbf{Thesis.} Mansfield lets us feel Miss Brill's Sunday communion from inside her own exclamations, then breaks it with a single overheard remark, revealing a sense of belonging that is vivid and sincere yet sustained entirely by her own imagination, so fragile that when it collapses she cannot even admit her own grief.

\textbf{Body 1: the play.} The free indirect exclamations (``Oh, how fascinating it was!'') put us inside her delight, and the stage metaphor converts loneliness into membership: ``Even she had a part.'' Yet the metaphor exposes its own weakness. She assures herself that ``No doubt somebody would have noticed'' her absence, the hedge of someone who is not sure; and her indifference to the invalid gentleman, whom she ``mightn't have noticed'' if he had died, is exactly the indifference the garden will soon show her.

\textbf{Body 2: the music.} The second paragraph rises toward a chorus in which ``all the whole company'' will sing, but Mansfield plants a ``faint chill'' that Miss Brill keeps renaming (``not sadness---no, not sadness''). The repetition is denial, and the ironic coda, ``though what they understood she didn't know,'' marks the gap between her feeling of understanding and any shared meaning.

\textbf{Body 3: the fall.} She casts the couple as ``hero and heroine,'' so their dialogue arrives as the play's real script, and in it she is ``that stupid old thing.'' Mansfield gives her no reply and no thoughts; the narration goes silent. The final paragraph's quiet details (the skipped honey-cake, the ``room like a cupboard'') shrink her world back to its true size, and the ``crying'' she hears from the fur's box is her own grief, displaced onto the one companion she has.

\textbf{Why it earns Row C.} The essay holds the tension throughout: the belonging is both genuine as feeling and false as fact, and the narration's shift from her voice to silence enacts the shift.
\end{modelessay}

\end{document}
